\begin{apartats}

\apartat[1,25]{0,75}
Deriva les funcions següents:
\begin{graella}{2}
  \sa f(x)=\left(1-x^2\right)^3 & \sa g(x)=2^{x^2}
\end{graella}

\begin{solucio}
i) $f'(x)=3\left(1-x^2\right)^2\cdot(-2x)=-6x\left(1-x^2\right)^2$.\\
ii) $g'(x)=2^{x^2}\ln2\cdot2x=2x\ln2\cdot2^{x^2}$.
\end{solucio}

\apartat[1,25]{1}
Deriva les funcions següents:
\begin{graella}{2}
  \sa h(x)=\ln(\sin x) & \sa k(x)=e^{\sqrt{x}}
\end{graella}

\begin{solucio}
i) $h'(x)=\dfrac{\cos x}{\sin x}$, per als $x$ on $\sin x>0$.\\
ii) $k'(x)=e^{\sqrt{x}}\cdot\dfrac{1}{2\sqrt{x}}=\dfrac{e^{\sqrt{x}}}{2\sqrt{x}}$, per a $x>0$.
\end{solucio}

\begin{nomesllarg}
\apartat{0,75}
Deriva la funció
\[
m(x)=\sqrt{1+\ln x}
\]
i calcula $m'(1)$.

\begin{solucio}
$m'(x)=\dfrac{1}{2\sqrt{1+\ln x}}\cdot\dfrac1x=\dfrac{1}{2x\sqrt{1+\ln x}}$.\\
En $x=1$: $\ln1=0$, i per tant $m'(1)=\dfrac{1}{2\cdot1\cdot1}=\dfrac12$.
\end{solucio}
\end{nomesllarg}

\end{apartats}
